% =============================================================
% บทที่ 5 - บทสรุป
% แก้ไขเนื้อหาในบทนี้ได้ที่ไฟล์ chapter5.tex
% =============================================================

% -----------------------------------------------------------
\section{สรุปผลการดำเนินโครงงาน}
% -----------------------------------------------------------

โครงงานนี้พัฒนาท่อส่งข้อมูลเชิงโมดูลาร์ตามกระบวนการวิทยาการข้อมูล ตั้งแต่การรวบรวมข้อมูล 5 ลีกใหญ่ของยุโรปจาก Understat, Transfermarkt และ Football-Data.co.uk การทำความสะอาดและตรวจสอบข้อมูล การสร้างตัวแปรก่อนแข่ง การพัฒนาแบบจำลองทำนายผล H/D/A ไปจนถึงการทดสอบย้อนหลังทางการเงิน โดยตั้งคำถามหลักว่า ``ผู้เดิมพันที่ใช้ข้อมูลสาธารณะจะเอาชนะตลาดเดิมพันฟุตบอลได้หรือไม่'' และออกแบบการทดสอบให้ตอบคำถามนี้อย่างตรงไปตรงมา ได้แก่ ทำนายทุกฤดูกาลจากฤดูกาลก่อนหน้าเท่านั้น ลงทะเบียนกลยุทธ์ไว้ล่วงหน้าก่อนเห็นผล และรายงานทุกผลลัพธ์พร้อมการปรับแก้การทดสอบหลายครั้ง ผลการดำเนินงานสรุปเทียบกับวัตถุประสงค์ได้ดังตารางที่ \ref{tab:obj_result}

\begin{table}[htbp]
\centering
\caption{สรุปผลเทียบกับวัตถุประสงค์ของโครงงาน}
\label{tab:obj_result}
\begin{tabular}{|c|>{\raggedright\arraybackslash}p{4.2cm}|>{\raggedright\arraybackslash}p{6.6cm}|c|}
\hline
\textbf{ข้อ} & \textbf{วัตถุประสงค์} & \textbf{ผลที่ได้} & \textbf{สถานะ} \\
\hline
1 & พัฒนากระบวนการวิทยาการข้อมูลแบบครบวงจรจากหลายแหล่งข้อมูล & ท่อส่งข้อมูลโมดูลาร์รวบรวมข้อมูล 5 ลีก ฤดูกาล 2014--2015 ถึง 2026--2027 เท่าที่แข่งไปแล้ว (21,882 นัดในตารางตัวแปร และ 20,712 นัดพร้อมราคาต่อรอง) ตรวจสอบทุกขั้นด้วย Pandera และสร้างผลซ้ำได้จากรหัสคอมมิต & บรรลุ \\ \hline
2 & สร้างและวิเคราะห์ตัวแปรก่อนแข่ง ได้แก่ Elo, ความเหนื่อยล้า, ตารางแข่ง และฟอร์มผู้เล่น 11 ตัวจริง & ยืนยันความได้เปรียบของเจ้าบ้านบนข้อมูลพรีเมียร์ลีก 2000--2024 ($+0.362$ ประตู) และความสัมพันธ์ของ Elo ($r = 0.441$) การคัดเลือกตัวแปรได้ Elo, ฟอร์มผู้เล่นตัวจริง, จำนวนนัดใน 14 วัน และนาทีลงเล่นก่อนหน้า ส่วนดัชนีความเหนื่อยล้าไม่ช่วยเมื่อมี Elo อยู่แล้ว & บรรลุ \\ \hline
3 & เปรียบเทียบความน่าจะเป็นของแบบจำลองกับราคาตลาดเพื่อค้นหา Value Bet & Logistic Regression ดีที่สุด (Log Loss $0.9855$, Accuracy $52.41\%$) แต่ด้อยกว่าราคาเจ้ามือ ($0.9691$) และการเพิ่มฟอร์มทีมเข้าไปในราคาตลาดไม่ทำให้ดีขึ้น ($+0.0008$) & บรรลุ \\ \hline
4 & ประเมินผลกำไรผ่านการทดสอบย้อนหลังด้วยกลยุทธ์การบริหารเงินทุนหลายรูปแบบ & ทดสอบ 18 กลยุทธ์ 42 รูปแบบใน 5 ลีก และแผนวางเงิน 5 แบบ การเดิมพันตามแบบจำลองขาดทุน $4$--$5.5\%$ ไม่มีกลยุทธ์ใดมีนัยสำคัญหลังปรับแก้ มีเพียง Line Shopping ที่มีกำไรคาดหมายเป็นบวก แต่เก็บได้ยากในทางปฏิบัติ & บรรลุ \\
\hline
\end{tabular}
\end{table}

โครงงานบรรลุวัตถุประสงค์ทั้ง 4 ข้อในแง่ของการพัฒนาระบบและการตอบคำถามวิจัย แต่คำตอบที่ได้เป็นเชิงลบ คือ \textbf{แบบจำลองที่สร้างจากข้อมูลสาธารณะไม่สามารถทำนายได้ดีกว่าตลาด และไม่มีกลยุทธ์ใดสร้างกำไรได้อย่างมีนัยสำคัญ} ซึ่งเป็นผลที่มีคุณค่าในตัวเอง เนื่องจากแสดงด้วยหลักฐานที่สร้างซ้ำได้ว่าตลาดเดิมพันของ 5 ลีกใหญ่มีประสิทธิภาพสูง และกำไรที่ปรากฏในการทดสอบย้อนหลังส่วนใหญ่อธิบายได้ด้วยความบังเอิญหรือการเลือกผลภายหลัง

\subsection{บทเรียนที่ได้จากโครงงาน}

\begin{enumerate}[label=\arabic*)]
    \item \textbf{ตลาดเป็นเกณฑ์เปรียบเทียบที่ถูกต้อง:} การวัดแบบจำลองเทียบกับการทายจากสัดส่วนผล (Class Prior) ทำให้ดูเหมือนแบบจำลองดีมาก แต่เมื่อเทียบกับราคาเจ้ามือจะเห็นว่ายังห่างอยู่ $0.0164$ Log Loss
    \item \textbf{แบบจำลองที่ซับซ้อนกว่าไม่ได้ดีกว่าเสมอ:} Logistic Regression ดีกว่า Gradient Boosting บนข้อมูลชุดนี้ เพราะสัญญาณส่วนใหญ่เป็นเชิงเส้นและมีขนาดเล็กเมื่อเทียบกับความแปรปรวนของผลการแข่งขัน
    \item \textbf{กำไรคาดหมาย ณ ราคาปิด น่าเชื่อถือกว่า ROI:} ROI ของฟุตบอลมีความแปรปรวนสูงจนต้องใช้เดิมพันหลายหมื่นครั้งจึงจะแยกความได้เปรียบเล็ก ๆ ออกจากโชคได้ ส่วนกำไรคาดหมาย ณ ราคาปิดให้ผลที่เสถียรกว่ามาก
    \item \textbf{การนับจำนวนครั้งที่ลองมีความสำคัญ:} เมื่อทดสอบ 252 ช่อง ย่อมมีบางช่องที่ ROI เป็นบวกโดยบังเอิญ การลงทะเบียนกลยุทธ์ล่วงหน้าและการปรับแก้แบบ Benjamini--Hochberg ช่วยป้องกันการสรุปผิด
    \item \textbf{การบริหารเงินทุนไม่สร้างความได้เปรียบ:} Martingale และ Kelly เต็มพาเงินทุนหมดได้ง่าย ส่วน Fractional Kelly ลดความเสี่ยงได้ แต่ไม่มีแผนใดเปลี่ยนค่าคาดหมายติดลบให้เป็นบวกได้
\end{enumerate}

% -----------------------------------------------------------
\section{ข้อจำกัดของระบบ}
% -----------------------------------------------------------

\subsection{ข้อจำกัดด้านการทำงาน}

\begin{itemize}
    \item \textbf{ความพร้อมของข้อมูลไลน์อัป (Lineup Availability):} รายชื่อผู้เล่น 11 ตัวจริงต้องรอการประกาศอย่างเป็นทางการล่วงหน้าก่อนแข่งประมาณ 60 นาที ทำให้แบบจำลองไม่สามารถประเมินผลในระยะยาวล่วงหน้าหลายวันได้
    \item \textbf{ความผันผวนของราคาอัตราต่อรอง (Odds Movement):} ระบบยังไม่ได้บันทึกความเคลื่อนไหวของราคาต่อรองแบบเรียลไทม์ (Live Odds) แต่ใช้อัตราต่อรองย้อนหลังคงที่จาก Football-Data.co.uk ซึ่งมีเพียงราคาเปิดและราคาปิด
    \item \textbf{ขอบเขตของตลาด:} ทดสอบเฉพาะ 5 ลีกใหญ่ และตลาด 1X2, Asian Handicap และ Over/Under 2.5 ลีกที่เล็กกว่าหรือตลาดอื่นอาจตั้งราคาได้แม่นน้อยกว่า แต่ก็เป็นตลาดที่เจ้ามือจำกัดบัญชีเร็วที่สุดเช่นกัน
    \item \textbf{ข้อมูล Transfermarkt สิ้นสุดที่ฤดูกาล 2025--2026:} ชุดข้อมูล transfermarkt-datasets หยุดอัปเดตในเดือนกรกฎาคม 2026 นัดของฤดูกาล 2026--2027 จึงไม่มีจำนวนนัดใน 14 วันจริง (ไม่ได้อยู่ในชุดทดสอบ แต่การทำนายนัดที่กำลังจะมาถึงต้องใช้ค่ามัธยฐานแทน) หากจะใช้งานต่อต้องหาแหล่งปฏิทินการแข่งขันทุกรายการแหล่งใหม่
    \item \textbf{ข้อมูลของแบบจำลองที่เทียบกับตลาด:} แบบจำลอง ``ราคาเปิด + ฟอร์มทีม'' ใช้ประตูและการยิงเข้ากรอบจาก Football-Data.co.uk ยังไม่ได้รวมฟอร์มผู้เล่นตัวจริงจาก Understat เข้าไป
    \item \textbf{สมมติฐานของการจำลอง:} การทดสอบที่ราคาดีที่สุดสมมติว่ามีบัญชีกับเจ้ามือทุกรายที่ Football-Data.co.uk บันทึกไว้ และราคาปิดเป็นกรณีที่ดีที่สุดของจังหวะเวลา ส่วนการจำลองโอกาสหมดตัวสุ่มเดิมพันทีละครั้ง จึงไม่รวมช่วงแพ้ติดกันภายในฤดูกาลและประเมินความเสี่ยงต่ำกว่าความจริง
\end{itemize}

\subsection{ข้อจำกัดด้านฮาร์ดแวร์และซอฟต์แวร์}

\begin{itemize}
    \item \textbf{หน่วยความจำและการประมวลผล:} การคำนวณตัวแปรระดับผู้เล่น (Player-level features) ย้อนหลัง 12 ฤดูกาลครอบคลุม 5 ลีกยุโรป ต้องใช้หน่วยความจำ RAM ขั้นต่ำ 8 GB และระบบประมวลผลแบบ Multi-core เพื่อให้รัน Data Pipeline ได้มีประสิทธิภาพ
    \item \textbf{ข้อจำกัดการเข้าถึง API แหล่งข้อมูล:} การดึงข้อมูลสถิติจากเว็บไซต์ภายนอกขึ้นอยู่กับโครงสร้าง HTML และขีดจำกัดการรับส่งข้อมูล (Rate Limit) ของเว็บไซต์ต้นทาง
\end{itemize}

% -----------------------------------------------------------
\section{ปัญหาอุปสรรค และแนวทางแก้ไข}
% -----------------------------------------------------------

ปัญหาที่พบระหว่างดำเนินงานและแนวทางแก้ไขแสดงดังตารางที่ \ref{tab:problems}

\begin{table}[htbp]
\centering
\caption{ปัญหาอุปสรรคและแนวทางแก้ไขในการดำเนินโครงงาน}
\label{tab:problems}
\begin{tabular}{|c|p{5cm}|p{6.5cm}|}
\hline
\textbf{ลำดับ} & \textbf{ปัญหาที่พบ} & \textbf{แนวทางแก้ไข} \\
\hline
1 & ชื่อสโมสรและนักเตะใน Understat, Transfermarkt และ Football-Data ไม่ตรงกัน & จัดทำระบบ Alias Table เพื่อทำ Entity Resolution แปลงชื่อให้เป็นมาตรฐานเดียวกันก่อนเชื่อมโยงข้อมูล \\ \hline
2 & รูปแบบและชนิดข้อมูลเกิดความคลาดเคลื่อนระหว่างขั้นตอนใน Pipeline & กำหนดสเปกข้อมูลด้วย Pandera Schema Validation ตรวจสอบโครงสร้างข้อมูลทุกสเตปก่อนส่งต่อไปฝึกโมเดล \\ \hline
3 & ปัญหาข้อมูลอนาคตรั่วไหล (Data Leakage) หากใช้สถิติหลังจบการแข่งขัน & ออกแบบระบบให้รองรับเฉพาะข้อมูลก่อนเริ่มแข่ง (Pre-match Data) และใช้วิธี Temporal Split ในการแบ่งชุดข้อมูลทดสอบ \\
\hline
\end{tabular}
\end{table}

% -----------------------------------------------------------
\section{ข้อเสนอแนะในการพัฒนาต่อไป}
% -----------------------------------------------------------

\subsection{การปรับปรุงฟังก์ชันการทำงาน}

\begin{itemize}
    \item \textbf{การเชื่อมต่อราคาเรียลไทม์ (Live Odds Integration):} พัฒนาระบบเชื่อมต่อ API กับเจ้ามือรับพนันเพื่อดึงราคาต่อรองล่าสุดก่อนเตะแบบอัตโนมัติ ทำให้คำนวณค่า $+EV$ ได้แม่นยำยิ่งขึ้น
    \item \textbf{การประมวลผลข่าวสารและความเคลื่อนไหว (Sentiment Analysis):} เพิ่มมอดูลวิเคราะห์ข่าวสารการบาดเจ็บ สภาพอากาศ หรือสัมภาษณ์ผู้จัดการทีมด้วยเทคนิค Natural Language Processing (NLP) เพื่อเสริมมิติการทำนาย
\end{itemize}

\subsection{แนวทางการวิจัยต่อยอด}

\begin{itemize}
    \item \textbf{การใช้ Deep Learning และ Graph Neural Networks (GNN):} ประยุกต์ใช้แบบจำลองโครงข่ายประสาทเทียมเพื่อเรียนรู้ความสัมพันธ์ระหว่างผู้เล่นในทีมแบบเชิงพื้นที่และปฏิสัมพันธ์บนสนาม
    \item \textbf{การขยายแบบจำลองสู่ตลาดการเดิมพันอื่น ๆ:} โครงงานนี้ทดสอบตลาดจำนวนประตูรวม (Over/Under 2.5) และเอเชียนแฮนดิแคป (Asian Handicap) เฉพาะกลยุทธ์ Line Shopping ที่อิงราคาของ Pinnacle ขั้นต่อไปคือใช้แบบจำลองจำนวนประตู Poisson ทำนายความน่าจะเป็นของตลาดเหล่านี้โดยตรง และทดสอบกับลีกรองที่ราคาอาจแม่นน้อยกว่า
\end{itemize}